% Source: https://artofproblemsolving.com/community/c3h2650413_sigma_mathematics_competition_smc_problems
% Stable identity: aops-topic-2650413-q14
Put $O$ at the origin and $D=\left(3\sqrt{13},0\right)$. For a point $P=10(\cos\phi,\sin\phi)$ of $\omega$, the law of cosines gives $DP^2=117+100-60\sqrt{13}\cos\phi$.
\begin{itemize}
\item $DB=1$: $\cos\phi_B=\dfrac{216}{60\sqrt{13}}=\dfrac{18}{5\sqrt{13}}$, so $\sin\phi_B=\pm\dfrac{1}{5\sqrt{13}}$ and $B=\dfrac{(36,\pm2)}{\sqrt{13}}$.
\item $DA=17$: $\cos\phi_A=-\dfrac{72}{60\sqrt{13}}=-\dfrac{6}{5\sqrt{13}}$, so $\sin\phi_A=\pm\dfrac{17}{5\sqrt{13}}$ and $A=\dfrac{(-12,\pm34)}{\sqrt{13}}$.
\item The tangent at $C$ is $\{X: X\cdot(\cos\phi_C,\sin\phi_C)=10\}$, at distance $\left|3\sqrt{13}\cos\phi_C-10\right|$ from $D$. This must equal $DB=1$. Since $3\sqrt{13}<11$, the only possibility is $3\sqrt{13}\cos\phi_C=9$, so $\cos\phi_C=\dfrac{3}{\sqrt{13}}$; reflecting in line $OD$ if necessary, $C=\dfrac{(30,20)}{\sqrt{13}}$.
\end{itemize}
The data do not fix the signs for $A$ and $B$. With all coordinates scaled by $\sqrt{13}$, the area is $\dfrac{1}{26}\left|\det\left(B-A,\,C-A\right)\right|$:
\begin{itemize}
\item $A=(-12,34)$, $B=(36,2)$: $\left|48\cdot(-14)-(-32)\cdot42\right|=672$, area $\dfrac{336}{13}$;
\item $A=(-12,-34)$, $B=(36,2)$: $\left|48\cdot54-36\cdot42\right|=1080$, area $\dfrac{540}{13}$;
\item $A=(-12,34)$, $B=(36,-2)$: $\left|48\cdot(-14)-(-36)\cdot42\right|=840$, area $\dfrac{420}{13}$;
\item $A=(-12,-34)$, $B=(36,-2)$: $\left|48\cdot54-32\cdot42\right|=1248$, area $48$.
\end{itemize}
The greatest possible area is $\boxed{48}$, when $A$ and $B$ lie on the opposite side of line $OD$ from $C$ (sketch; $T$ is the foot of the perpendicular from $D$ to the tangent, $DT=DB=1$).

\begin{asy}
usepackage("amsmath");
size(5.5cm);
defaultpen(fontsize(10pt));
real r=10;pair O=(0,0),D=(3*sqrt(13),0);
pair C=(30,20)/sqrt(13),B=(36,-2)/sqrt(13),A=(-12,-34)/sqrt(13);
pair u=unit(C),T=D+(r-dot(D,u))*u;
draw(circle(O,r),linewidth(0.8));fill(A--B--C--cycle,lightgray);draw(A--B--C--cycle,linewidth(0.8));
draw(C+9*(-u.y,u.x)--C-8*(-u.y,u.x));draw(D--T,dashed);draw(D--B,dashed);draw(O--D,dotted);draw(O--C,dotted);
dot(O);dot(D);dot(T);dot(B);
label("$O$",O,W);label("$A$",A,SW);label("$B$",B,SW);label("$C$",C,NE);label("$D$",D,SE,fontsize(9pt));label("$T$",T,E,fontsize(9pt));
\end{asy}
